\begin{abcliste}
\abc
Indifferent ist das Gleichgewicht nur, wenn die Achse durch den \emph{Schwerpunkt} geht. Die Mitte des Bretts ist aber nicht der Schwerpunkt des ganzen Körpers: Das Blei verschiebt ihn nach rechts. Die Gewichtskraft (von Brett und Blei zusammen) hat bezüglich D also einen Hebelarm, ausser wenn S senkrecht unter oder über D liegt.

\abc
Die Achse ist der Nullpunkt. Der Schwerpunkt des Bretts liegt in D, das Blei im Abstand $L/2$:
\begin{align}
 x_\mathrm{S} &= \xSF \\
   &= \frac{\mPb\cdot\frac{\LB}{2}}{\mB+\mPb} \\
   &= \xS \approx \result{\xSP{-2}{2}}
\end{align}
rechts von D, auf dem Brett zum Blei hin.

\abc
Das Gewicht des Bretts greift in D an und hat keinen Hebelarm; es bleibt das Drehmoment des Bleis (gleichwertig: die ganze Gewichtskraft $(m_\mathrm{B}+m_\mathrm{Pb})\,g$ mit dem Hebelarm $x_\mathrm{S}$):
\begin{align}
 M &= \MhF = \mPb\cdot\g\cdot\frac{\LB}{2} = \Mh \approx \result{\MhP{0}{2}}
\end{align}
Waagrecht ist das Brett also nicht im Gleichgewicht: Es dreht sich im Uhrzeigersinn, das Blei nach unten.

\abc
Das Brett pendelt und kommt zur Ruhe, wenn S senkrecht unter D liegt: \result{\text{Brett senkrecht, Blei unten}}. Das ist ein \result{\text{stabiles Gleichgewicht}}: Dreht man das Brett ein wenig, hebt sich S, und die Gewichtskraft dreht es zurück.

Eine zweite Gleichgewichtslage gibt es mit dem Blei senkrecht oben (S senkrecht über D). Sie ist labil: Schon eine kleine Drehung lässt das Brett umkippen.
\end{abcliste}
